\documentclass[11pt,a4paper]{article}

% Packages
\usepackage[utf8]{inputenc}
\usepackage[T1]{fontenc}
\usepackage{lmodern}
\usepackage[margin=1in]{geometry}
\usepackage{graphicx}
\usepackage{hyperref}
\usepackage{listings}
\usepackage{xcolor}
\usepackage{booktabs}
\usepackage{longtable}
\usepackage{float}
% enumitem not available, using standard enumerate/itemize
\usepackage{amsmath}
\usepackage{amssymb}
\usepackage{fancyhdr}
% titlesec and tocloft not available in basic installation
% \usepackage{titlesec}
% \usepackage{tocloft}
\setlength{\parskip}{0.5em}
\setlength{\parindent}{0em}

% Hyperref setup
\hypersetup{
    colorlinks=true,
    linkcolor=blue,
    filecolor=magenta,
    urlcolor=cyan,
    pdftitle={GLIDER Technical Documentation},
    pdfauthor={GLIDER Development Team},
}

% Code listing style
\definecolor{codegreen}{rgb}{0,0.6,0}
\definecolor{codegray}{rgb}{0.5,0.5,0.5}
\definecolor{codepurple}{rgb}{0.58,0,0.82}
\definecolor{backcolour}{rgb}{0.95,0.95,0.92}

\lstdefinestyle{mystyle}{
    backgroundcolor=\color{backcolour},
    commentstyle=\color{codegreen},
    keywordstyle=\color{magenta},
    numberstyle=\tiny\color{codegray},
    stringstyle=\color{codepurple},
    basicstyle=\ttfamily\footnotesize,
    breakatwhitespace=false,
    breaklines=true,
    captionpos=b,
    keepspaces=true,
    numbers=left,
    numbersep=5pt,
    showspaces=false,
    showstringspaces=false,
    showtabs=false,
    tabsize=2,
    frame=single
}
\lstset{style=mystyle}

% Header/Footer
\pagestyle{fancy}
\fancyhf{}
\rhead{GLIDER Technical Documentation}
\lhead{Version 0.3.0}
\cfoot{\thepage}

% Using default title formatting (titlesec not available)

\begin{document}

% Title Page
\begin{titlepage}
    \centering
    \vspace*{2cm}

    {\Huge\bfseries GLIDER}\\[0.5cm]
    {\Large\textit{General Laboratory Interface for Design,\\Experimentation, and Recording}}\\[2cm]

    {\LARGE Technical Documentation}\\[0.5cm]
    {\large Version 0.3.0}\\[3cm]

    {\large A Comprehensive Reference for the\\Visual Laboratory Automation Framework}\\[2cm]

    \vfill

    {\large\textbf{Authors:}\\GLIDER Development Team}\\[1cm]

    {\large February 2026}

\end{titlepage}

% Abstract
\begin{abstract}
\noindent
GLIDER (General Laboratory Interface for Design, Experimentation, and Recording) is a modular experimental orchestration platform designed for laboratory hardware control through visual flow-based programming. This technical documentation provides a comprehensive reference for researchers, developers, and laboratory technicians working with the GLIDER framework. The system enables non-programmers to design and execute complex experiments through a node-based visual interface while providing extensibility for advanced users. Key features include hardware abstraction for Arduino and Raspberry Pi platforms, integrated computer vision with YOLO-based object tracking, real-time data recording, and an AI-powered assistant for experimental guidance. This document covers the complete system architecture, all major components, installation procedures, and usage patterns derived from extensive codebase analysis.
\end{abstract}

\newpage
\tableofcontents
\newpage

% ============================================================================
\section{Introduction}
% ============================================================================

\subsection{Overview}

GLIDER represents a paradigm shift in laboratory automation, providing researchers with a powerful yet accessible tool for designing, executing, and analyzing experiments. Unlike traditional laboratory software that requires programming expertise, GLIDER employs a visual node-based programming paradigm similar to game engine blueprint systems, making complex experimental protocols accessible to researchers without coding experience.

\subsection{Key Features}

\begin{itemize}    \item \textbf{Visual Flow Programming:} Node-based interface for experiment design
    \item \textbf{Hardware Abstraction:} Unified interface for Arduino and Raspberry Pi
    \item \textbf{Computer Vision Integration:} YOLO-based object detection and tracking
    \item \textbf{Multi-Camera Support:} Synchronized recording from multiple sources
    \item \textbf{Behavioral Analysis:} Automated detection of freeze, dart, and mobility states
    \item \textbf{Zone-Based Tracking:} Spatial analysis with customizable regions of interest
    \item \textbf{Data Recording:} Comprehensive CSV logging with metadata
    \item \textbf{AI Assistant:} LLM-powered guidance for experiment design
    \item \textbf{Cross-Platform:} Windows, macOS, Linux, and Raspberry Pi support
    \item \textbf{Dual UI Modes:} Desktop IDE and touch-optimized runner interface
\end{itemize}

\subsection{Target Applications}

GLIDER is designed for behavioral and biological research laboratories requiring:
\begin{itemize}    \item Animal behavior tracking and analysis
    \item Operant conditioning protocols
    \item Optogenetics and stimulation experiments
    \item Environmental chamber control
    \item Multi-modal data acquisition
    \item Real-time experimental automation
\end{itemize}

\subsection{Document Organization}

This documentation is organized into the following sections:
\begin{enumerate}    \item \textbf{System Architecture:} High-level design and component relationships
    \item \textbf{Core Components:} Central orchestration and state management
    \item \textbf{Hardware Abstraction Layer:} Board and device drivers
    \item \textbf{Visual Programming System:} Node types and flow execution
    \item \textbf{Computer Vision:} Camera management and tracking
    \item \textbf{Graphical User Interface:} Desktop and runner modes
    \item \textbf{Data Recording:} CSV logging and video capture
    \item \textbf{Serialization:} Experiment file format
    \item \textbf{Plugin System:} Extensibility framework
    \item \textbf{Installation \& Usage:} Setup and operation guides
    \item \textbf{API Reference:} Detailed component documentation
\end{enumerate}

% ============================================================================
\section{System Architecture}
% ============================================================================

\subsection{Architectural Overview}

GLIDER employs a layered Model-View-Controller (MVC) architecture with clear separation of concerns. The system is organized into distinct layers that communicate through well-defined interfaces and callback mechanisms.

\begin{figure}[H]
\centering
\begin{verbatim}
+-------------------------------------------------------------+
|                    GUI Layer (PyQt6)                         |
|  - Main Window, Panels, Dialogs, Node Graph View, Runner    |
+------------------------------+------------------------------+
                               |
+------------------------------v------------------------------+
|              GliderCore (Central Orchestrator)               |
|  - Coordinates all subsystems, manages lifecycle and state  |
+------------------------------+------------------------------+
                               |
        +----------------------+----------------------+
        |                      |                      |
+-------v--------+   +---------v--------+   +--------v--------+
|  Flow Engine   |   | Hardware Manager |   |  Vision System  |
|  (Execution)   |   |  (HAL Control)   |   |   (CV/Tracking) |
+-------+--------+   +---------+--------+   +--------+--------+
        |                      |                      |
+-------v--------+   +---------v--------+   +--------v--------+
| Node Registry  |   |  Board Drivers   |   | Camera Manager  |
| Flow Graph     |   |  Device Registry |   | CV Processor    |
+----------------+   +------------------+   +-----------------+
\end{verbatim}
\caption{GLIDER System Architecture}
\end{figure}

\subsection{Component Relationships}

The following diagram illustrates the relationships between major components:

\begin{verbatim}
GliderCore (Orchestrator)
|
+-- ExperimentSession (Data Model)
|   +-- SessionMetadata
|   +-- HardwareConfig (boards, devices)
|   +-- FlowConfig (nodes, connections)
|   +-- CameraConfig
|   +-- ZoneConfig
|   +-- Subject[] (animal/experiment subjects)
|
+-- HardwareManager (Device Lifecycle)
|   +-- Board Registry
|   |   +-- TelemetrixBoard (Arduino)
|   |   +-- PiGPIOBoard (Raspberry Pi)
|   |   +-- MockBoard (Testing)
|   +-- Device Registry
|   |   +-- DigitalOutputDevice
|   |   +-- DigitalInputDevice
|   |   +-- AnalogInputDevice
|   |   +-- PWMOutputDevice
|   |   +-- ServoDevice
|   +-- PinManager (conflict detection)
|
+-- FlowEngine (Graph Execution)
|   +-- Node Registry
|   +-- Connection Graph
|   +-- Task Management
|   +-- CustomDeviceRunner
|   +-- FlowFunctionRunner
|
+-- Vision System
|   +-- CameraManager / MultiCameraManager
|   +-- CVProcessor
|   +-- BehaviorAnalyzer
|   +-- ZoneTracker
|   +-- VideoRecorder / MultiVideoRecorder
|   +-- TrackingLogger
|
+-- DataRecorder (CSV Logging)
+-- AgentController (AI Assistant)
\end{verbatim}

\subsection{Design Patterns}

GLIDER employs several well-established software design patterns:

\begin{table}[H]
\centering
\begin{tabular}{@{}lp{10cm}@{}}
\toprule
\textbf{Pattern} & \textbf{Application} \\
\midrule
Observer & Callback-based state change notifications throughout the system \\
Factory & Device, node, and board creation through registry lookup \\
Strategy & Interchangeable CV backends (background subtraction, YOLO) \\
Command & Undo/redo stack for GUI operations \\
Abstract Factory & Hardware abstraction layer with pluggable drivers \\
Singleton & GliderCore as central application coordinator \\
State Machine & Session state management (IDLE, RUNNING, PAUSED, etc.) \\
\bottomrule
\end{tabular}
\caption{Design Patterns in GLIDER}
\end{table}

\subsection{Threading Model}

GLIDER uses a multi-threaded architecture to maintain UI responsiveness:

\begin{itemize}    \item \textbf{Main Thread:} PyQt6 event loop, all UI updates
    \item \textbf{Camera Thread:} Background frame capture via CameraManager
    \item \textbf{CV Thread:} Dedicated QThread for computer vision processing
    \item \textbf{Hardware Thread:} Telemetrix runs in background thread for Arduino
    \item \textbf{Recording Thread:} Video encoding in separate thread
\end{itemize}

Thread safety is maintained through:
\begin{itemize}    \item Signal marshaling for cross-thread communication
    \item Frame copying to avoid race conditions
    \item QThread for managed background processing
    \item No direct UI access from background threads
\end{itemize}

% ============================================================================
\section{Core Components}
% ============================================================================

\subsection{GliderCore}

The \texttt{GliderCore} class serves as the central orchestrator for the entire application, managing component lifecycle, state transitions, and inter-component communication.

\subsubsection{Responsibilities}

\begin{itemize}    \item Initialize and manage the asyncio event loop
    \item Create and coordinate HardwareManager, FlowEngine, and vision components
    \item Manage ExperimentSession (the data model)
    \item Load plugins and register built-in nodes
    \item Orchestrate recording and video streaming
    \item Handle state callbacks and error propagation
\end{itemize}

\subsubsection{Key Properties}

\begin{lstlisting}[language=Python]
class GliderCore:
    session: ExperimentSession          # Current experiment state
    hardware_manager: HardwareManager   # Device and board management
    flow_engine: FlowEngine             # Logic graph execution
    camera_manager: CameraManager       # Single camera interface
    multi_camera_manager: MultiCameraManager
    data_recorder: DataRecorder         # CSV data logging
    video_recorder: VideoRecorder       # Video capture
    tracking_logger: TrackingLogger     # Animal tracking data
    cv_processor: CVProcessor           # Real-time vision processing
\end{lstlisting}

\subsubsection{Lifecycle Methods}

\begin{lstlisting}[language=Python]
async def initialize():
    """Initialize all subsystems"""

async def start_experiment():
    """Begin experiment execution"""
    # 1. Setup hardware
    # 2. Connect boards
    # 3. Setup flow graph
    # 4. Start recording
    # 5. Start flow engine

async def stop_experiment():
    """Stop experiment and cleanup"""
    # 1. Stop flow engine
    # 2. Stop all recorders
    # 3. Set devices to safe state

async def shutdown():
    """Graceful application shutdown"""
\end{lstlisting}

\subsection{ExperimentSession}

The \texttt{ExperimentSession} class represents the complete state of an experiment, serving as the data model that can be serialized to and from \texttt{.glider} files.

\subsubsection{Session State Machine}

\begin{verbatim}
IDLE --> INITIALIZING --> READY --> RUNNING <--> PAUSED
  ^                         |          |
  |                         v          v
  +-------- STOPPING <------+----------+
                |
                v
              ERROR
\end{verbatim}

\subsubsection{Data Structure}

\begin{lstlisting}[language=Python]
@dataclass
class ExperimentSession:
    # Metadata
    name: str
    description: str
    author: str
    protocol: str
    experiment_type: str
    experimenter: str
    lab: str
    project: str
    notes: str

    # Configuration
    hardware_config: HardwareConfig
    flow_config: FlowConfig
    camera_config: CameraConfig
    zone_config: ZoneConfig
    dashboard_config: DashboardConfig

    # Subjects
    subjects: List[Subject]
    active_subject_id: Optional[str]

    # State tracking
    state: SessionState
    is_dirty: bool
    file_path: Optional[Path]
\end{lstlisting}

\subsubsection{Subject Data Model}

\begin{lstlisting}[language=Python]
@dataclass
class Subject:
    subject_id: str          # User-defined identifier (e.g., "M001")
    age: str                 # Age in appropriate units
    sex: str                 # M/F/Unknown
    weight: str              # Weight with units
    strain: str              # Animal strain/genotype
    genotype: str            # Genetic modifications
    solution: str            # Drug/treatment solution
    concentration: str       # Solution concentration
    dose: str                # Administered dose
    route: str               # Administration route (IP/IV/PO/SC)
    group: str               # Treatment group assignment
    notes: str               # Additional observations
\end{lstlisting}

\subsection{FlowEngine}

The \texttt{FlowEngine} manages execution of the visual programming graph, supporting both data flow (reactive) and execution flow (imperative) paradigms.

\subsubsection{Execution Models}

\textbf{Data Flow (Reactive):}
\begin{itemize}    \item Values propagate automatically when inputs change
    \item Used for continuous sensor monitoring and data transformation
    \item Similar to spreadsheet formulas
\end{itemize}

\textbf{Execution Flow (Imperative):}
\begin{itemize}    \item Explicit sequential execution triggered by EXEC ports
    \item Used for experiment protocols with defined steps
    \item Similar to traditional programming
\end{itemize}

\subsubsection{Flow States}

\begin{verbatim}
STOPPED --> RUNNING <--> PAUSED --> ERROR
    ^           |
    +-----------+
\end{verbatim}

\subsubsection{Node Registry}

The FlowEngine maintains a registry of available node types:

\begin{lstlisting}[language=Python]
class FlowEngine:
    _node_registry: Dict[str, Type[GliderNode]]

    def register_node(self, node_type: str, node_class: Type):
        """Register a node type for use in flow graphs"""

    def create_node(self, node_id: str, node_type: str,
                    state: dict, device_id: str = None):
        """Instantiate a node from the registry"""
\end{lstlisting}

\subsection{HardwareManager}

The \texttt{HardwareManager} provides lifecycle management for boards and devices, including connection handling, pin allocation, and error recovery.

\subsubsection{Board Management}

\begin{lstlisting}[language=Python]
class HardwareManager:
    async def create_board(self, config: BoardConfig) -> BaseBoard:
        """Create and register a board from configuration"""

    async def connect_board(self, board_id: str):
        """Establish connection to hardware"""

    async def disconnect_board(self, board_id: str):
        """Gracefully disconnect board"""
\end{lstlisting}

\subsubsection{Device Management}

\begin{lstlisting}[language=Python]
    async def create_device(self, config: DeviceConfig) -> BaseDevice:
        """Create device with pin validation"""
        # 1. Validate pin availability
        # 2. Allocate pins via PinManager
        # 3. Create device instance
        # 4. Register in device registry

    async def initialize_device(self, device_id: str):
        """Configure pins and set initial state"""
\end{lstlisting}

\subsubsection{Pin Manager}

The \texttt{PinManager} prevents pin conflicts:

\begin{lstlisting}[language=Python]
class PinManager:
    def allocate_pin(self, pin: int, device_id: str,
                     mode: PinMode) -> bool:
        """Allocate pin to device, returns False if conflict"""

    def release_pin(self, pin: int, device_id: str):
        """Release pin allocation"""

    def get_pin_owner(self, pin: int) -> Optional[str]:
        """Get device ID that owns a pin"""
\end{lstlisting}

% ============================================================================
\section{Hardware Abstraction Layer}
% ============================================================================

\subsection{Overview}

The Hardware Abstraction Layer (HAL) provides a unified interface for interacting with different hardware platforms. This abstraction allows the same experiment to run on Arduino, Raspberry Pi, or simulated hardware without code changes.

\subsection{BaseBoard Interface}

\begin{lstlisting}[language=Python]
class BaseBoard(ABC):
    """Abstract base class for hardware board drivers"""

    @abstractmethod
    async def connect(self) -> bool:
        """Establish connection to hardware"""

    @abstractmethod
    async def disconnect(self):
        """Close hardware connection"""

    @abstractmethod
    async def set_pin_mode(self, pin: int, mode: PinMode,
                           pin_type: PinType):
        """Configure pin for specific operation"""

    @abstractmethod
    async def write_digital(self, pin: int, value: bool):
        """Set digital output state"""

    @abstractmethod
    async def read_digital(self, pin: int) -> bool:
        """Read digital input state"""

    @abstractmethod
    async def write_pwm(self, pin: int, value: int):
        """Set PWM duty cycle (0-255)"""

    @abstractmethod
    async def read_analog(self, pin: int) -> int:
        """Read analog input (0-1023)"""
\end{lstlisting}

\subsection{Board Implementations}

\subsubsection{TelemetrixBoard (Arduino)}

The \texttt{TelemetrixBoard} provides communication with Arduino boards via the Telemetrix protocol over USB serial.

\begin{table}[H]
\centering
\begin{tabular}{@{}ll@{}}
\toprule
\textbf{Feature} & \textbf{Specification} \\
\midrule
Protocol & Telemetrix-AIO \\
Connection & USB Serial \\
Supported Boards & Arduino Uno, Mega, Nano, etc. \\
Digital Pins & Board-dependent \\
Analog Pins & A0-A5 (Uno), A0-A15 (Mega) \\
PWM Pins & Board-dependent \\
Latency & ~50ms per operation \\
\bottomrule
\end{tabular}
\caption{TelemetrixBoard Specifications}
\end{table}

\subsubsection{PiGPIOBoard (Raspberry Pi)}

The \texttt{PiGPIOBoard} provides direct GPIO access on Raspberry Pi via the pigpio library.

\begin{table}[H]
\centering
\begin{tabular}{@{}ll@{}}
\toprule
\textbf{Feature} & \textbf{Specification} \\
\midrule
Library & pigpio / gpiozero \\
GPIO Pins & 28 (BCM numbering) \\
PWM & Hardware PWM on GPIO 12, 13, 18, 19 \\
Analog & Requires external ADC \\
Latency & <5ms (RPi 4+) \\
\bottomrule
\end{tabular}
\caption{PiGPIOBoard Specifications}
\end{table}

\subsubsection{MockBoard (Testing)}

The \texttt{MockBoard} provides an in-memory simulation for testing without physical hardware.

\subsection{BaseDevice Interface}

\begin{lstlisting}[language=Python]
class BaseDevice(ABC):
    """Abstract base class for semantic device wrappers"""

    @property
    def actions(self) -> Dict[str, Callable]:
        """Map of action names to async methods"""

    async def execute_action(self, action: str, *args, **kwargs):
        """Execute a named action"""

    def get_state(self) -> Any:
        """Get current device state"""

    async def initialize(self):
        """Configure pins and initial state"""

    async def shutdown(self):
        """Cleanup and safe state"""
\end{lstlisting}

\subsection{Built-in Device Types}

\begin{table}[H]
\centering
\begin{tabular}{@{}llp{7cm}@{}}
\toprule
\textbf{Device Type} & \textbf{Pins} & \textbf{Actions} \\
\midrule
DigitalOutputDevice & signal & on, off, toggle, set(value) \\
DigitalInputDevice & signal & read() returns bool \\
AnalogInputDevice & signal & read() returns 0-1023 \\
PWMOutputDevice & signal & set\_duty(0-255), set\_percent(0-100) \\
ServoDevice & signal & set\_angle(0-180) \\
StepperDevice & step, dir, enable & step(count), set\_speed(rpm) \\
\bottomrule
\end{tabular}
\caption{Built-in Device Types}
\end{table}

% ============================================================================
\section{Visual Programming System}
% ============================================================================

\subsection{Node Architecture}

GLIDER's visual programming system is built on the ryvencore library, extended with custom node types for laboratory applications.

\subsubsection{GliderNode Base Class}

\begin{lstlisting}[language=Python]
class GliderNode(ABC):
    """Base class for all GLIDER nodes"""

    # Port definitions
    inputs: List[PortDefinition]
    outputs: List[PortDefinition]

    # Node metadata
    node_type: str
    category: str
    visible_in_runner: bool = False

    # Device binding
    device_id: Optional[str]

    async def execute(self):
        """Execute node logic (ExecNode)"""

    def process(self, **inputs) -> Dict[str, Any]:
        """Process inputs and return outputs (DataNode)"""
\end{lstlisting}

\subsubsection{Node Categories}

\begin{enumerate}    \item \textbf{Experiment Nodes:} StartExperiment, EndExperiment, Delay
    \item \textbf{Control Nodes:} Loop, WaitForInput, Sequence, Timer
    \item \textbf{Hardware Nodes:} DigitalWrite, DigitalRead, AnalogRead, PWMWrite
    \item \textbf{Logic Nodes:} Add, Subtract, Multiply, Divide, MapRange, Clamp
    \item \textbf{Comparison Nodes:} Threshold, InRange
    \item \textbf{Interface Nodes:} Button, Slider, Label, Gauge, Chart
    \item \textbf{Vision Nodes:} ZoneInput
\end{enumerate}

\subsection{Node Catalog}

\subsubsection{Experiment Nodes}

\begin{table}[H]
\centering
\begin{tabular}{@{}lp{8cm}@{}}
\toprule
\textbf{Node} & \textbf{Description} \\
\midrule
StartExperiment & Entry point for experiment flow; triggers when experiment starts \\
EndExperiment & Terminates experiment execution; signals completion \\
Delay & Pauses execution for specified duration (seconds) \\
Output & Writes a value to a bound device \\
Input & Reads current value from a bound device \\
MotorGovernor & Controls stepper motor speed and direction \\
CustomDevice & Wrapper for user-defined device logic \\
\bottomrule
\end{tabular}
\caption{Experiment Nodes}
\end{table}

\subsubsection{Control Nodes}

\begin{table}[H]
\centering
\begin{tabular}{@{}lp{8cm}@{}}
\toprule
\textbf{Node} & \textbf{Description} \\
\midrule
Loop & Repeats connected flow N times \\
WaitForInput & Pauses until user interaction or signal \\
Sequence & Executes outputs in order with optional delays \\
Timer & Triggers output after duration elapses \\
Toggle & Alternates between two states on each trigger \\
\bottomrule
\end{tabular}
\caption{Control Nodes}
\end{table}

\subsubsection{Hardware Nodes}

\begin{table}[H]
\centering
\begin{tabular}{@{}lp{8cm}@{}}
\toprule
\textbf{Node} & \textbf{Description} \\
\midrule
DigitalWrite & Sets digital output HIGH or LOW \\
DigitalRead & Reads digital input state (bool) \\
AnalogRead & Reads analog input (0-1023) \\
PWMWrite & Sets PWM duty cycle (0-255) \\
DeviceAction & Executes named action on bound device \\
DeviceRead & Reads state from bound device \\
\bottomrule
\end{tabular}
\caption{Hardware Nodes}
\end{table}

\subsubsection{Logic Nodes}

\begin{table}[H]
\centering
\begin{tabular}{@{}lp{8cm}@{}}
\toprule
\textbf{Node} & \textbf{Description} \\
\midrule
Add & Outputs sum of two inputs \\
Subtract & Outputs difference of two inputs \\
Multiply & Outputs product of two inputs \\
Divide & Outputs quotient of two inputs \\
MapRange & Linearly maps input from one range to another \\
Clamp & Constrains value between min and max \\
Threshold & Outputs true if input exceeds threshold \\
InRange & Outputs true if input is within range \\
PID & Proportional-Integral-Derivative controller \\
\bottomrule
\end{tabular}
\caption{Logic Nodes}
\end{table}

\subsubsection{Interface Nodes}

\begin{table}[H]
\centering
\begin{tabular}{@{}lp{8cm}@{}}
\toprule
\textbf{Node} & \textbf{Description} \\
\midrule
Label & Displays text value in runner dashboard \\
Gauge & Circular meter display for numeric values \\
Chart & Time-series line chart \\
LEDIndicator & Visual on/off indicator \\
Button & Momentary push button for user input \\
ToggleSwitch & On/off switch control \\
Slider & Continuous value input via slider \\
NumericInput & Numeric entry with spinner buttons \\
\bottomrule
\end{tabular}
\caption{Interface Nodes}
\end{table}

\subsection{Flow Functions}

Flow Functions allow creating reusable experiment subroutines:

\begin{lstlisting}[language=Python]
# Define a flow function
StartFunction("blink_led")  # Entry point
  --> Delay(0.5)
  --> Output("on")
  --> Delay(0.5)
  --> Output("off")
  --> EndFunction()         # Exit point

# Call the function elsewhere
FunctionCall("blink_led")   # Invokes the subroutine
\end{lstlisting}

\subsection{Connection Types}

\subsubsection{Data Connections}

Data connections (blue) carry values between nodes:
\begin{itemize}    \item Propagate automatically when source changes
    \item Support type coercion where applicable
    \item Can have multiple destinations (fan-out)
\end{itemize}

\subsubsection{Execution Connections}

Execution connections (white) control flow sequencing:
\begin{itemize}    \item Trigger target node when source completes
    \item Sequential execution guarantee
    \item Single destination per output (no fan-out)
\end{itemize}

% ============================================================================
\section{Computer Vision System}
% ============================================================================

\subsection{Overview}

GLIDER integrates comprehensive computer vision capabilities for automated behavioral analysis, including object detection, multi-object tracking, and spatial analysis.

\subsection{CameraManager}

The \texttt{CameraManager} provides platform-agnostic camera access with support for various camera types.

\subsubsection{Supported Cameras}

\begin{table}[H]
\centering
\begin{tabular}{@{}lll@{}}
\toprule
\textbf{Camera Type} & \textbf{Platform} & \textbf{Backend} \\
\midrule
USB Webcam & All & OpenCV (DirectShow/V4L2) \\
Scientific (Y800) & All & FFmpeg subprocess \\
CSI Camera & Raspberry Pi & picamera2 \\
Miniscope & Linux & v4l2-ctl + OpenCV \\
\bottomrule
\end{tabular}
\caption{Supported Camera Types}
\end{table}

\subsubsection{Camera Settings}

\begin{lstlisting}[language=Python]
@dataclass
class CameraSettings:
    resolution: Tuple[int, int] = (640, 480)
    fps: int = 30
    exposure: int = -6
    brightness: int = 128
    contrast: int = 128
    saturation: int = 128
    auto_focus: bool = True
    auto_exposure: bool = True
    miniscope_mode: bool = False
    pixel_format: Optional[str] = None  # e.g., "Y800"
    connection_timeout: float = 5.0
    buffer_size: int = 2
\end{lstlisting}

\subsection{CVProcessor}

The \texttt{CVProcessor} performs real-time image analysis with configurable detection backends.

\subsubsection{Detection Backends}

\begin{table}[H]
\centering
\begin{tabular}{@{}lp{7cm}l@{}}
\toprule
\textbf{Backend} & \textbf{Description} & \textbf{Speed} \\
\midrule
BACKGROUND\_SUBTRACTION & MOG2 background modeling with morphological filtering & Fast \\
YOLO\_V8 & Single-frame object detection & Medium \\
YOLO\_BYTETRACK & Multi-object tracking with persistent IDs & Slower \\
MOTION\_ONLY & Simple frame difference thresholding & Fastest \\
\bottomrule
\end{tabular}
\caption{CV Detection Backends}
\end{table}

\subsubsection{Processing Pipeline}

\begin{verbatim}
Frame Input
    |
    v
+-------------------+
| Detection Backend |  (Background sub / YOLO / Motion)
+-------------------+
    |
    v
+-------------------+
| Contour/BBox      |  (Extract detected regions)
| Extraction        |
+-------------------+
    |
    v
+-------------------+
| Object Tracker    |  (Centroid matching, ID assignment)
+-------------------+
    |
    v
+-------------------+
| Behavior Analyzer |  (Freeze/dart/mobility detection)
+-------------------+
    |
    v
+-------------------+
| Zone Checker      |  (ROI occupancy)
+-------------------+
    |
    v
Output: TrackedObjects[], ZoneStates[], BehaviorStates
\end{verbatim}

\subsection{Object Tracking}

\subsubsection{TrackedObject Data Structure}

\begin{lstlisting}[language=Python]
@dataclass
class TrackedObject:
    track_id: int           # Persistent identifier
    class_name: str         # Object class (e.g., "mouse")
    bbox: Tuple[int,int,int,int]  # x, y, width, height
    confidence: float       # Detection confidence (0-1)
    centroid: Tuple[int,int]      # Center point
    age: int                # Frames since first detection
    disappeared: int        # Frames since last seen
    behavioral_state: str   # freeze/dart/immobile/unknown
    velocity: float         # Pixels per frame
\end{lstlisting}

\subsubsection{Centroid Tracking Algorithm}

\begin{enumerate}    \item Detect objects in current frame (bounding boxes)
    \item Extract centroids from detections
    \item Compute distance matrix between existing tracks and new detections
    \item Perform greedy minimum-distance assignment
    \item Mark unmatched tracks as ``disappeared''
    \item Create new tracks for unmatched detections
    \item Deregister tracks exceeding \texttt{max\_disappeared} threshold
\end{enumerate}

\subsection{Behavior Analysis}

The \texttt{BehaviorAnalyzer} classifies animal behavior from tracking data.

\subsubsection{Behavioral States}

\begin{table}[H]
\centering
\begin{tabular}{@{}lp{8cm}@{}}
\toprule
\textbf{State} & \textbf{Criteria} \\
\midrule
FREEZE & velocity $<$ freeze\_threshold for $\geq$ freeze\_duration frames \\
IMMOBILE & velocity $<$ immobile\_threshold \\
DART & velocity $>$ dart\_threshold (sudden movement) \\
EXPLORING & Normal movement within thresholds \\
UNKNOWN & Insufficient tracking data \\
\bottomrule
\end{tabular}
\caption{Behavioral State Classification}
\end{table}

\subsection{Zone System}

\subsubsection{Zone Types}

\begin{itemize}    \item \textbf{Rectangle:} Defined by two corner points
    \item \textbf{Circle:} Defined by center and radius point
    \item \textbf{Polygon:} Arbitrary vertices (3+ points)
\end{itemize}

\subsubsection{Zone Tracking}

\begin{lstlisting}[language=Python]
@dataclass
class ZoneState:
    zone_id: str
    zone_name: str
    occupied: bool
    entry_time: Optional[float]
    exit_time: Optional[float]
    time_inside: float
\end{lstlisting}

\subsubsection{Spatial Queries}

\begin{lstlisting}[language=Python]
class Zone:
    def contains_point(self, norm_x: float, norm_y: float) -> bool:
        """Check if normalized coordinates are within zone"""

    def contains_point_pixels(self, px: int, py: int,
                              width: int, height: int) -> bool:
        """Check if pixel coordinates are within zone"""
\end{lstlisting}

\subsection{Video Recording}

\subsubsection{VideoRecorder}

\begin{lstlisting}[language=Python]
class VideoRecorder:
    def __init__(self, output_path: Path,
                 resolution: Tuple[int,int],
                 fps: int = 30,
                 annotated: bool = False):
        """Initialize video recording"""

    def record_frame(self, frame: np.ndarray,
                     tracks: List[TrackedObject] = None):
        """Write frame with optional tracking overlay"""

    def stop(self):
        """Finalize video file"""
\end{lstlisting}

\subsubsection{MultiVideoRecorder}

Coordinates synchronized recording from multiple cameras with configurable layouts.

\subsection{Calibration}

\subsubsection{Camera Calibration}

\begin{lstlisting}[language=Python]
class CameraCalibration:
    def add_line(self, start: Tuple[int,int],
                 end: Tuple[int,int],
                 length: float, unit: LengthUnit,
                 resolution: Tuple[int,int]):
        """Add calibration reference line"""

    @property
    def pixels_per_mm(self) -> float:
        """Get calibrated conversion factor"""

    def pixels_to_mm(self, pixels: float) -> float:
        """Convert pixel distance to millimeters"""

    def mm_to_pixels(self, mm: float) -> float:
        """Convert millimeters to pixels"""
\end{lstlisting}

\subsubsection{Supported Units}

\begin{itemize}    \item Millimeters (mm)
    \item Centimeters (cm)
    \item Meters (m)
    \item Inches (in)
    \item Feet (ft)
\end{itemize}

% ============================================================================
\section{Graphical User Interface}
% ============================================================================

\subsection{Overview}

GLIDER provides two user interface modes optimized for different use cases:

\begin{itemize}    \item \textbf{Builder Mode (Desktop):} Full IDE for experiment design
    \item \textbf{Runner Mode (Touch):} Optimized dashboard for experiment execution
\end{itemize}

\subsection{MainWindow}

The \texttt{MainWindow} serves as the application shell, managing view switching and global operations.

\subsubsection{Structure}

\begin{verbatim}
MainWindow
+-- MenuBar
|   +-- File (New, Open, Save, Export)
|   +-- Edit (Undo, Redo, Preferences)
|   +-- View (Builder, Runner, Fullscreen)
|   +-- Hardware (Connect, Disconnect, Configure)
|   +-- Experiment (Start, Stop, Pause)
|   +-- Help (Documentation, About)
|
+-- ToolBar
|   +-- Quick actions (New, Save, Run, Stop)
|
+-- StackedWidget
|   +-- Index 0: Builder View
|   +-- Index 1: Runner View
|
+-- StatusBar
    +-- Session state, connection status
\end{verbatim}

\subsection{Builder Mode Components}

\subsubsection{Node Graph Editor}

The visual scripting canvas built on Qt Graphics View framework:

\begin{itemize}    \item \textbf{NodeGraphScene:} Infinite canvas with grid background
    \item \textbf{NodeGraphView:} Pan/zoom interaction, drag-and-drop
    \item \textbf{NodeItem:} Visual representation of nodes
    \item \textbf{PortItem:} Input/output connection points
    \item \textbf{ConnectionItem:} Bezier curve connections
\end{itemize}

\textbf{Interaction:}
\begin{itemize}    \item Pan: Middle mouse or Ctrl+Left drag
    \item Zoom: Mouse wheel (0.1x to 3.0x)
    \item Connect: Click-drag between ports
    \item Select: Click node, Shift+click for multi-select
    \item Delete: Delete key or context menu
\end{itemize}

\subsubsection{Experiment Panel}

Manages experiment metadata and subjects:

\begin{verbatim}
+----------------------------------+
| Experiment Info (Collapsible)    |
| - Name, Protocol, Type           |
| - Experimenter, Lab, Project     |
| - Notes (text area)              |
+----------------------------------+
| Subjects (Table)                 |
| [ID | Name | Group | Sex | ...]  |
| [Add] [Edit] [Remove] [Active]   |
+----------------------------------+
| Active Subject Display (Green)   |
| Active: M001 - Mouse (Group A)   |
+----------------------------------+
\end{verbatim}

\subsubsection{Camera Panel}

Live camera preview with CV controls:

\begin{itemize}    \item Camera preview with overlay support
    \item Resolution and FPS display
    \item CV toggle and backend selection
    \item Recording controls
    \item Zone visualization
    \item Calibration line overlay
\end{itemize}

\subsubsection{Agent Panel}

AI assistant chat interface:

\begin{itemize}    \item Message history with role-based styling
    \item Quick prompt buttons
    \item Streaming response display
    \item Action confirmation dialogs
\end{itemize}

\subsection{Runner Mode}

Touch-optimized interface for 480x800 portrait displays (Raspberry Pi):

\begin{verbatim}
+----------------------------------+ 50px
| [Experiment Name] [Timer] [Gear] | Header
+----------------------------------+
| [Recording Indicator]            | 28px
+----------------------------------+
| Device Status Cards (Scrollable) | Flex
| +------------------------------+ |
| | Device 1: ON                 | |
| +------------------------------+ |
| | Device 2: 127                | |
| +------------------------------+ |
+----------------------------------+
| [  START  ]  [  STOP  ]         | 60px
| [    EMERGENCY STOP    ]        | 60px
+----------------------------------+ 160px
\end{verbatim}

\subsubsection{Touch Widgets}

Large, high-contrast widgets for touch interaction:

\begin{table}[H]
\centering
\begin{tabular}{@{}lll@{}}
\toprule
\textbf{Widget} & \textbf{Min Height} & \textbf{Purpose} \\
\midrule
TouchLabel & 80px & Display text values \\
TouchButton & 60px & Momentary press \\
TouchToggle & 80px & On/off switch \\
TouchSlider & 80px & Continuous value \\
TouchGauge & 120px & Circular meter \\
TouchChart & 200px & Time series \\
TouchLED & 80px & Status indicator \\
TouchNumericInput & 80px & Number entry \\
\bottomrule
\end{tabular}
\caption{Touch Widget Specifications}
\end{table}

\subsection{Dialogs}

\begin{table}[H]
\centering
\begin{tabular}{@{}lp{8cm}@{}}
\toprule
\textbf{Dialog} & \textbf{Purpose} \\
\midrule
CameraSettingsDialog & Resolution, FPS, format, CV settings \\
ZoneDialog & Draw and edit zones on camera preview \\
CalibrationDialog & Pixel-to-real-world calibration \\
FlowFunctionDialog & Custom Python function editor \\
AnalysisDialog & Post-experiment data analysis \\
ExperimentDialog & Full experiment setup wizard \\
AgentSettingsDialog & AI agent configuration \\
SubjectDialog & Subject/animal metadata \\
CustomDeviceDialog & Define custom device types \\
\bottomrule
\end{tabular}
\caption{Dialog Reference}
\end{table}

\subsection{Command Pattern (Undo/Redo)}

GLIDER implements full undo/redo support:

\begin{lstlisting}[language=Python]
class CreateNodeCommand(Command):
    def execute(self):
        # Create node in graph

    def undo(self):
        # Remove node from graph

# Command types:
# - CreateNodeCommand
# - DeleteNodeCommand
# - MoveNodeCommand
# - PropertyChangeCommand
# - CreateConnectionCommand
# - DeleteConnectionCommand
\end{lstlisting}

\subsection{View Mode Detection}

Automatic UI mode selection based on display:

\begin{lstlisting}[language=Python]
def detect_view_mode(screen_size: Tuple[int,int]) -> ViewMode:
    width, height = screen_size

    # Raspberry Pi touchscreen (480x800 or 800x480)
    if (width, height) in [(480, 800), (800, 480)]:
        return ViewMode.RUNNER

    # Small displays
    if width <= 480:
        return ViewMode.RUNNER

    return ViewMode.DESKTOP
\end{lstlisting}

% ============================================================================
\section{Data Recording and Analysis}
% ============================================================================

\subsection{DataRecorder}

The \texttt{DataRecorder} generates timestamped CSV files with comprehensive metadata.

\subsubsection{Output Format}

\begin{lstlisting}
# GLIDER Experiment Data
# Experiment Name: My Experiment
# Start Time: 2026-02-03T10:30:45.123456
# Sample Interval (s): 0.1
# Protocol: Social Interaction
# Experimenter: John Smith
# Lab: Behavioral Neuroscience Lab
# Subject: M001
# ----- Boards -----
# arduino_1: Arduino Uno on /dev/ttyUSB0
# ----- Devices -----
# led_1: DigitalOutput on pin 13
# sensor_1: AnalogInput on pin A0

timestamp,elapsed_ms,led_1,sensor_1,zone:center
2026-02-03T10:30:45.123,0.0,0,512,0
2026-02-03T10:30:45.223,100.0,1,515,0
2026-02-03T10:30:45.323,200.0,1,518,1
\end{lstlisting}

\subsubsection{Configuration}

\begin{lstlisting}[language=Python]
class DataRecorder:
    def __init__(self,
                 output_dir: Path,
                 sample_interval_ms: int = 100,
                 include_zones: bool = True):
        """Initialize recorder with configuration"""
\end{lstlisting}

\subsection{TrackingLogger}

Logs detailed tracking data for behavioral analysis:

\begin{lstlisting}
timestamp,frame,track_id,class,confidence,x,y,w,h,cx,cy,velocity,state
2026-02-03T10:30:45.123,0,1,mouse,0.95,100,150,50,60,125,180,0.0,unknown
2026-02-03T10:30:45.156,1,1,mouse,0.93,102,152,50,60,127,182,2.8,exploring
2026-02-03T10:30:45.189,2,1,mouse,0.94,103,153,50,60,128,183,1.4,freeze
\end{lstlisting}

\subsection{Video Recording}

\subsubsection{Standard Recording}

Raw video capture at configured resolution and FPS.

\subsubsection{Annotated Recording}

Video with tracking overlays:
\begin{itemize}    \item Bounding boxes around tracked objects
    \item Track ID labels
    \item Centroid markers
    \item Trajectory trails
    \item Zone boundaries
    \item Timestamp overlay
\end{itemize}

\subsection{Export Formats}

\begin{table}[H]
\centering
\begin{tabular}{@{}lll@{}}
\toprule
\textbf{Data Type} & \textbf{Format} & \textbf{Extension} \\
\midrule
Device States & CSV & .csv \\
Tracking Data & CSV & \_tracking.csv \\
Video & MP4 (H.264) & .mp4 \\
Annotated Video & MP4 (H.264) & \_annotated.mp4 \\
Experiment Config & JSON & .glider \\
\bottomrule
\end{tabular}
\caption{Data Export Formats}
\end{table}

% ============================================================================
\section{Serialization System}
% ============================================================================

\subsection{Experiment File Format}

GLIDER experiments are saved as JSON files with the \texttt{.glider} extension.

\subsubsection{Schema Structure}

\begin{lstlisting}[language=json]
{
  "metadata": {
    "name": "Experiment Name",
    "version": "1.0.0",
    "description": "...",
    "author": "...",
    "experiment_type": "...",
    "protocol": "...",
    "experimenter": "...",
    "lab": "...",
    "project": "...",
    "notes": "...",
    "created": "2026-02-03T10:00:00",
    "modified": "2026-02-03T12:00:00",
    "subjects": [...]
  },
  "hardware": {
    "boards": [...],
    "devices": [...]
  },
  "flow": {
    "nodes": [...],
    "connections": [...]
  },
  "camera": {...},
  "zones": {...},
  "dashboard": {...},
  "custom_devices": [...],
  "flow_functions": [...]
}
\end{lstlisting}

\subsubsection{Board Configuration}

\begin{lstlisting}[language=json]
{
  "id": "arduino_1",
  "driver_type": "arduino",
  "port": "/dev/ttyUSB0",
  "board_type": "uno",
  "auto_reconnect": true,
  "settings": {}
}
\end{lstlisting}

\subsubsection{Device Configuration}

\begin{lstlisting}[language=json]
{
  "id": "led_1",
  "board_id": "arduino_1",
  "device_type": "DigitalOutput",
  "name": "Status LED",
  "pins": {"signal": 13},
  "settings": {"inverted": false}
}
\end{lstlisting}

\subsubsection{Node Configuration}

\begin{lstlisting}[language=json]
{
  "id": "node_001",
  "type": "Delay",
  "position": {"x": 200, "y": 150},
  "properties": {"duration": 1.0},
  "device_id": null
}
\end{lstlisting}

\subsubsection{Connection Configuration}

\begin{lstlisting}[language=json]
{
  "id": "conn_001",
  "from_node": "node_001",
  "from_port": 0,
  "to_node": "node_002",
  "to_port": 0,
  "connection_type": "exec"
}
\end{lstlisting}

\subsection{Schema Validation}

GLIDER validates experiment files against JSON schemas:

\begin{lstlisting}[language=Python]
from glider.serialization.schema import validate_experiment

errors = validate_experiment(data)
if errors:
    for error in errors:
        print(f"Validation error at {error.path}: {error.message}")
\end{lstlisting}

% ============================================================================
\section{Plugin System}
% ============================================================================

\subsection{Overview}

GLIDER's plugin system enables extensibility through:
\begin{itemize}    \item Custom hardware drivers
    \item Custom device types
    \item Custom node types
\end{itemize}

\subsection{Plugin Discovery}

Plugins are discovered from:
\begin{enumerate}    \item Python entry points (\texttt{glider.plugins})
    \item Local plugin directory (\texttt{\~{}/.glider/plugins/})
\end{enumerate}

\subsection{Creating a Hardware Driver}

\begin{lstlisting}[language=Python]
from glider.hal.base_board import BaseBoard

class MyCustomBoard(BaseBoard):
    """Custom hardware driver"""

    @classmethod
    def get_driver_name(cls) -> str:
        return "my_custom_board"

    async def connect(self) -> bool:
        # Implement connection logic
        pass

    async def write_digital(self, pin: int, value: bool):
        # Implement digital write
        pass

    # ... implement other abstract methods

# Register in pyproject.toml:
# [project.entry-points."glider.boards"]
# my_custom_board = "my_package:MyCustomBoard"
\end{lstlisting}

\subsection{Creating a Custom Node}

\begin{lstlisting}[language=Python]
from glider.nodes.base_node import GliderNode, DataNode

class MyCustomNode(DataNode):
    """Custom processing node"""

    node_type = "MyCustomNode"
    category = "Custom"

    inputs = [
        PortDefinition("input", PortType.DATA, float)
    ]

    outputs = [
        PortDefinition("output", PortType.DATA, float)
    ]

    def process(self, input: float) -> Dict[str, Any]:
        result = input * 2  # Custom logic
        return {"output": result}

# Register in pyproject.toml:
# [project.entry-points."glider.nodes"]
# my_custom_node = "my_package:MyCustomNode"
\end{lstlisting}

% ============================================================================
\section{Installation and Usage}
% ============================================================================

\subsection{System Requirements}

\begin{table}[H]
\centering
\begin{tabular}{@{}ll@{}}
\toprule
\textbf{Component} & \textbf{Requirement} \\
\midrule
Python & 3.9, 3.10, 3.11, 3.12, or 3.13 (not 3.14+) \\
Operating System & Windows 10+, macOS 10.15+, Linux, Raspberry Pi OS \\
Memory & 4GB RAM minimum, 8GB recommended \\
Storage & 500MB for installation \\
Display & 1024x768 minimum (Desktop), 480x800 (Runner) \\
\bottomrule
\end{tabular}
\caption{System Requirements}
\end{table}

\subsection{Installation}

\subsubsection{Desktop (Windows/macOS/Linux)}

\begin{lstlisting}[language=bash]
# Create virtual environment
python -m venv venv
source venv/bin/activate  # Linux/macOS
# or: venv\Scripts\Activate  # Windows

# Install with desktop dependencies
pip install "glider[pc]"

# Launch application
glider
\end{lstlisting}

\subsubsection{Raspberry Pi}

\begin{lstlisting}[language=bash]
# Install system dependencies
sudo apt install python3-pyqt6 python3-venv

# Create virtual environment with system packages
python -m venv --system-site-packages venv
source venv/bin/activate

# Install GLIDER
pip install -e .

# Launch application
glider
\end{lstlisting}

\subsubsection{Development Installation}

\begin{lstlisting}[language=bash]
# Clone repository
git clone https://github.com/user/glider.git
cd glider

# Create virtual environment
python -m venv venv
source venv/bin/activate

# Install with development dependencies
pip install -e ".[dev]"

# Run tests
pytest
\end{lstlisting}

\subsection{Command Line Options}

\begin{lstlisting}[language=bash]
glider [OPTIONS]

Options:
  --builder      Force builder (desktop) mode
  --runner       Force runner (touch) mode
  --file PATH    Open experiment file on startup
  --debug        Enable debug logging
  --help         Show help message
\end{lstlisting}

\subsection{Quick Start Guide}

\subsubsection{Step 1: Create New Experiment}

\begin{enumerate}    \item Launch GLIDER
    \item File $\rightarrow$ New Experiment
    \item Enter experiment metadata (name, protocol, etc.)
    \item Click Create
\end{enumerate}

\subsubsection{Step 2: Configure Hardware}

\begin{enumerate}    \item Hardware $\rightarrow$ Add Board
    \item Select board type (Arduino/Raspberry Pi)
    \item Configure port/connection settings
    \item Hardware $\rightarrow$ Add Device
    \item Select device type and configure pins
\end{enumerate}

\subsubsection{Step 3: Design Flow}

\begin{enumerate}    \item Drag nodes from library to canvas
    \item Connect nodes by clicking and dragging between ports
    \item Configure node properties in inspector
    \item Bind hardware nodes to devices
\end{enumerate}

\subsubsection{Step 4: Configure Camera (Optional)}

\begin{enumerate}    \item View $\rightarrow$ Camera Panel
    \item Select camera from dropdown
    \item Configure resolution and FPS
    \item Enable CV processing if needed
    \item Draw zones for spatial analysis
\end{enumerate}

\subsubsection{Step 5: Run Experiment}

\begin{enumerate}    \item File $\rightarrow$ Save
    \item Experiment $\rightarrow$ Start
    \item Monitor in Runner view
    \item Experiment $\rightarrow$ Stop when complete
\end{enumerate}

\subsection{Troubleshooting}

\begin{table}[H]
\centering
\begin{tabular}{@{}p{5cm}p{8cm}@{}}
\toprule
\textbf{Issue} & \textbf{Solution} \\
\midrule
Arduino not detected & Check USB connection, install drivers, verify port permissions \\
Camera not working & Check device permissions, try different backend, verify resolution \\
CV processing slow & Reduce resolution, use background subtraction backend \\
UI unresponsive & Check for hardware communication issues, reduce sample rate \\
Import errors & Verify all dependencies installed, check Python version \\
\bottomrule
\end{tabular}
\caption{Common Issues and Solutions}
\end{table}

% ============================================================================
\section{API Reference}
% ============================================================================

\subsection{Core Module}

\subsubsection{glider.core.glider\_core}

\begin{lstlisting}[language=Python]
class GliderCore:
    """Central application orchestrator"""

    session: ExperimentSession
    hardware_manager: HardwareManager
    flow_engine: FlowEngine

    async def initialize()
    async def start_experiment()
    async def stop_experiment()
    async def shutdown()
\end{lstlisting}

\subsubsection{glider.core.experiment\_session}

\begin{lstlisting}[language=Python]
class ExperimentSession:
    """Experiment state model"""

    # Properties
    name: str
    state: SessionState
    hardware_config: HardwareConfig
    flow_config: FlowConfig

    # Methods
    def to_dict() -> dict
    @classmethod
    def from_dict(data: dict) -> ExperimentSession
    def save(path: Path)
    @classmethod
    def load(path: Path) -> ExperimentSession
\end{lstlisting}

\subsection{HAL Module}

\subsubsection{glider.hal.base\_board}

\begin{lstlisting}[language=Python]
class BaseBoard(ABC):
    """Abstract hardware driver interface"""

    @abstractmethod
    async def connect() -> bool
    @abstractmethod
    async def disconnect()
    @abstractmethod
    async def set_pin_mode(pin: int, mode: PinMode, type: PinType)
    @abstractmethod
    async def write_digital(pin: int, value: bool)
    @abstractmethod
    async def read_digital(pin: int) -> bool
    @abstractmethod
    async def write_pwm(pin: int, value: int)
    @abstractmethod
    async def read_analog(pin: int) -> int
\end{lstlisting}

\subsection{Vision Module}

\subsubsection{glider.vision.camera\_manager}

\begin{lstlisting}[language=Python]
class CameraManager:
    """Camera capture management"""

    def enumerate_cameras() -> List[CameraInfo]
    def open(camera_index: int, settings: CameraSettings)
    def start_streaming()
    def stop_streaming()
    def get_frame() -> Optional[np.ndarray]
    def add_frame_callback(callback: Callable)
\end{lstlisting}

\subsubsection{glider.vision.cv\_processor}

\begin{lstlisting}[language=Python]
class CVProcessor:
    """Computer vision processing"""

    def set_backend(backend: CVBackend)
    def process_frame(frame: np.ndarray) -> CVResult
    def get_tracked_objects() -> List[TrackedObject]
    def get_zone_states() -> List[ZoneState]
\end{lstlisting}

\subsection{Nodes Module}

\subsubsection{glider.nodes.base\_node}

\begin{lstlisting}[language=Python]
class GliderNode(ABC):
    """Base node class"""

    node_type: str
    category: str
    inputs: List[PortDefinition]
    outputs: List[PortDefinition]

    def get_state() -> dict
    def set_state(state: dict)
    def on_output_update(callback: Callable)

class DataNode(GliderNode):
    """Reactive data processing node"""
    def process(**inputs) -> Dict[str, Any]

class ExecNode(GliderNode):
    """Imperative execution node"""
    async def execute()
\end{lstlisting}

% ============================================================================
\section{Testing}
% ============================================================================

\subsection{Test Organization}

\begin{verbatim}
tests/
+-- conftest.py              # Shared pytest fixtures
+-- unit/                    # Unit tests
|   +-- core/
|   |   +-- test_config.py
|   |   +-- test_experiment_session.py
|   |   +-- test_flow_engine.py
|   |   +-- test_types.py
|   +-- hal/
|   |   +-- test_pin_manager.py
|   |   +-- test_telemetrix_board.py
|   +-- nodes/
|   |   +-- test_base_node.py
|   +-- serialization/
|   |   +-- test_schema.py
|   +-- vision/
|       +-- test_calibration.py
|       +-- test_zones.py
+-- integration/
    +-- test_experiment_workflow.py
\end{verbatim}

\subsection{Running Tests}

\begin{lstlisting}[language=bash]
# Run all tests
pytest

# Run with coverage
pytest --cov=glider

# Run specific module
pytest tests/unit/core/

# Run with verbose output
pytest -v

# Run only fast tests (skip slow integration tests)
pytest -m "not slow"
\end{lstlisting}

\subsection{Test Fixtures}

Key fixtures available in \texttt{conftest.py}:

\begin{lstlisting}[language=Python]
@pytest.fixture
def mock_board() -> MockBoard:
    """Provides mock hardware board for testing"""

@pytest.fixture
def mock_hardware_manager() -> HardwareManager:
    """Provides configured hardware manager"""

@pytest.fixture
def sample_session() -> ExperimentSession:
    """Provides sample experiment session"""

@pytest.fixture
def mock_camera_manager() -> CameraManager:
    """Provides mock camera for testing"""
\end{lstlisting}

% ============================================================================
\section{Performance Characteristics}
% ============================================================================

\begin{table}[H]
\centering
\begin{tabular}{@{}lll@{}}
\toprule
\textbf{Component} & \textbf{Typical Latency} & \textbf{Notes} \\
\midrule
Hardware write (Arduino) & 50-100ms & Serial communication overhead \\
Hardware write (RPi) & <5ms & Direct GPIO access \\
Camera frame capture & 16-33ms & Depends on FPS setting \\
CV processing (background sub) & 10-30ms & Resolution dependent \\
CV processing (YOLO) & 30-100ms & GPU acceleration recommended \\
Data recording & <1ms & Async CSV write \\
Node execution & 0-5ms & Depends on device latency \\
\bottomrule
\end{tabular}
\caption{Performance Characteristics}
\end{table}

% ============================================================================
\section{Safety Considerations}
% ============================================================================

\subsection{Hardware Safety}

\begin{itemize}    \item \textbf{Emergency Stop:} Sets all outputs to safe (LOW) state immediately
    \item \textbf{Pin Conflict Detection:} Prevents multiple devices using same pin
    \item \textbf{Graceful Shutdown:} Proper cleanup on application exit
    \item \textbf{Auto-Reconnection:} Handles temporary connection loss
\end{itemize}

\subsection{Data Safety}

\begin{itemize}    \item \textbf{Dirty State Tracking:} Warns before closing unsaved experiments
    \item \textbf{Async Recording:} Prevents data loss from UI blocking
    \item \textbf{Final Flush:} Ensures all data written on stop
    \item \textbf{Metadata Embedding:} Complete context in output files
\end{itemize}

\subsection{Experiment Safety}

\begin{itemize}    \item \textbf{Flow Validation:} Checks for required nodes before execution
    \item \textbf{State Callbacks:} Immediate notification of errors
    \item \textbf{Subject Tracking:} Complete metadata for reproducibility
\end{itemize}

% ============================================================================
\section{Conclusion}
% ============================================================================

GLIDER provides a comprehensive platform for laboratory automation that balances accessibility with power. The visual programming paradigm enables researchers without coding experience to design complex experiments, while the extensible architecture allows advanced users to add custom functionality.

Key strengths include:
\begin{itemize}    \item Intuitive visual experiment design
    \item Flexible hardware abstraction supporting multiple platforms
    \item Integrated computer vision for behavioral analysis
    \item Comprehensive data recording and export
    \item Touch-optimized interface for embedded deployment
    \item AI-powered assistance for experiment guidance
\end{itemize}

Future development directions may include:
\begin{itemize}    \item Additional hardware driver support
    \item Enhanced behavioral analysis algorithms
    \item Cloud integration for remote monitoring
    \item Expanded node library
    \item Machine learning integration
\end{itemize}

% ============================================================================
\appendix
\section{Appendix: File Structure}
% ============================================================================

\begin{verbatim}
/src/glider/
+-- __init__.py
+-- __main__.py              # Application entry point
+-- core/
|   +-- glider_core.py       # Central orchestrator
|   +-- experiment_session.py # Data model
|   +-- flow_engine.py       # Graph execution
|   +-- hardware_manager.py  # Device lifecycle
|   +-- data_recorder.py     # CSV logging
|   +-- flow_function.py     # Subroutine definitions
|   +-- custom_device.py     # User device logic
|   +-- config.py            # Configuration
|   +-- types.py             # Type definitions
+-- hal/
|   +-- base_board.py        # Driver interface
|   +-- base_device.py       # Device interface
|   +-- pin_manager.py       # Pin allocation
|   +-- mock_board.py        # Testing driver
|   +-- boards/
|       +-- telemetrix_board.py  # Arduino
|       +-- pi_gpio_board.py     # Raspberry Pi
+-- vision/
|   +-- camera_manager.py
|   +-- multi_camera_manager.py
|   +-- cv_processor.py
|   +-- zones.py
|   +-- behavior_analyzer.py
|   +-- calibration.py
|   +-- video_recorder.py
|   +-- tracking_logger.py
+-- nodes/
|   +-- base_node.py
|   +-- experiment_nodes.py
|   +-- control_nodes.py
|   +-- flow_function_nodes.py
|   +-- hardware/
|   +-- logic/
|   +-- interface/
|   +-- vision/
+-- serialization/
|   +-- serializer.py
|   +-- schema.py
+-- gui/
|   +-- main_window.py
|   +-- panels/
|   +-- dialogs/
|   +-- node_graph/
|   +-- runner/
|   +-- widgets/
|   +-- styles/
+-- agent/
|   +-- agent_controller.py
|   +-- llm_backend.py
|   +-- toolkit.py
|   +-- tools/
+-- plugins/
    +-- plugin_manager.py
\end{verbatim}

% ============================================================================
\section{Appendix: Dependencies}
% ============================================================================

\subsection{Core Dependencies}

\begin{table}[H]
\centering
\begin{tabular}{@{}lll@{}}
\toprule
\textbf{Package} & \textbf{Version} & \textbf{Purpose} \\
\midrule
qasync & $\geq$0.24.0 & Async Qt integration \\
telemetrix-aio & $\geq$1.0.0 & Arduino communication \\
pyserial & $\geq$3.5 & Serial port access \\
jsonschema & $\geq$4.0.0 & Schema validation \\
pyqtgraph & $\geq$0.13.0 & Scientific plotting \\
numpy & $\geq$1.21.0 & Numerical arrays \\
opencv-python & $\geq$4.8.0 & Computer vision \\
scipy & $\geq$1.10.0 & Scientific computing \\
httpx & $\geq$0.24.0 & HTTP client \\
ryvencore & $\geq$0.4.0 & Node-based flows \\
ultralytics & $\geq$8.0.0 & YOLO models \\
pandas & $\geq$2.0.0 & Data analysis \\
\bottomrule
\end{tabular}
\caption{Core Dependencies}
\end{table}

\subsection{Optional Dependencies}

\begin{table}[H]
\centering
\begin{tabular}{@{}lll@{}}
\toprule
\textbf{Extra} & \textbf{Packages} & \textbf{Purpose} \\
\midrule
pc & PyQt6 & Desktop GUI \\
rpi & gpiozero, lgpio, picamera2 & Raspberry Pi support \\
dev & pytest, black, ruff & Development tools \\
\bottomrule
\end{tabular}
\caption{Optional Dependencies}
\end{table}

% ============================================================================
\section{Appendix: Configuration Reference}
% ============================================================================

\subsection{Timing Configuration}

\begin{lstlisting}[language=Python]
@dataclass
class TimingConfig:
    device_refresh_interval_ms: int = 250
    elapsed_timer_interval_ms: int = 1000
    default_poll_interval: float = 0.1
    board_ready_timeout: float = 10.0
    board_operation_timeout: float = 5.0
    function_execution_timeout: float = 60.0
\end{lstlisting}

\subsection{UI Configuration}

\begin{lstlisting}[language=Python]
@dataclass
class UIConfig:
    min_window_width: int = 1024
    min_window_height: int = 768
    default_window_width: int = 1400
    default_window_height: int = 900
    max_undo_stack_size: int = 100
\end{lstlisting}

\subsection{Hardware Configuration}

\begin{lstlisting}[language=Python]
@dataclass
class HardwareConfig:
    adc_resolution: int = 1023
    adc_reference_voltage: float = 5.0
    pwm_min: int = 0
    pwm_max: int = 255
    servo_min_angle: int = 0
    servo_max_angle: int = 180
    input_poll_interval_ms: int = 100
\end{lstlisting}

\end{document}
